% GENERATED FILE. Edit canonical lessons, then run npm run generate:books.
% canonical-source-hash: fnv1a64:4525240bba921d60
% canonical-lessons: GE-C527-trockner, GE-C527-waescheleine, GE-C527-waescheklammer, GE-C527-buegeln, GE-C527-falten

\chapter{Wäsche waschen}
\label{ch:ge-wasche-waschen}
\hlchaptermodality{527}

\begin{chapteropening}
\textbf{By the end of this chapter:} \emph{I can talk about doing the laundry.}

Complete the last lesson of chapter 527: I can talk about doing the laundry.
\end{chapteropening}

\section[der Trockner]{der Trockner — tumble dryer}

\label{lesson:GE-C527-trockner}



\begin{quote}
Before the new one: say the German for to vacuum, then the German for cloth.
\end{quote}

\subsection*{You'll want to know: der Trockner}
\textbf{der Trockner} — \textquotedblleft{}tumble dryer\textquotedblright{}.

\begin{quote}
Die Handtücher kommen in den Trockner.
\end{quote}

\begin{quote}
\emph{The towels go in the dryer.}
\end{quote}

Plural: \emph{die Trockner}. Built on \emph{trocknen} — to dry.

\begin{cousinweb}
\emph{Trockner} is built on \emph{trocknen} \textquotedblleft{}to dry\textquotedblright{}, from \emph{trocken} \textquotedblleft{}dry\textquotedblright{}, which is inherited from Middle High German \emph{trucken} (EtymDB).
\end{cousinweb}

\begin{grammarlens}[title={What you've built}]
One more word for telling, explaining and describing at B1.
\end{grammarlens}

\subsection*{Your turn}
\begin{itemize}
\raggedright
  \item \emph{Say it:} \emph{der Trockner}
  \item \emph{Say it:} \emph{der Trockner}, once more
  \item \emph{Say it:} a sentence of your own with \emph{der Trockner}
  \item \emph{Recall:} say the German for to vacuum, then the German for cloth, then say \emph{der Trockner} again
\end{itemize}

\begin{tcolorbox}[breakable,colback=teal!4,colframe=teal!35!black,title={Before you move on}]
What does der Trockner mean? (\textquotedblleft{}Tumble dryer\textquotedblright{}.) Say it once more.
\end{tcolorbox}
\section[die Wäscheleine]{die Wäscheleine — clothes line}

\label{lesson:GE-C527-waescheleine}



\begin{quote}
Before the new one: say the German for cloth, then the German for tumble dryer.
\end{quote}

\subsection*{You'll want to know: die Wäscheleine}
\textbf{die Wäscheleine} — \textquotedblleft{}clothes line\textquotedblright{}.

\begin{quote}
Im Sommer hängen wir die Wäsche auf die Wäscheleine.
\end{quote}

\begin{quote}
\emph{In summer we hang the washing on the line.}
\end{quote}

Plural: \emph{die Wäscheleinen}.

\begin{cousinweb}
\emph{Wäscheleine} is built on \emph{Wäsche} \textquotedblleft{}washing\textquotedblright{} and \emph{Leine} \textquotedblleft{}line, rope\textquotedblright{}, inherited from Old High German \emph{līna}, from Proto-Germanic \emph{līnōn} (EtymDB).
\end{cousinweb}

\begin{grammarlens}[title={What you've built}]
One more word for telling, explaining and describing at B1.
\end{grammarlens}

\subsection*{Your turn}
\begin{itemize}
\raggedright
  \item \emph{Say it:} \emph{die Wäscheleine}
  \item \emph{Say it:} \emph{die Wäscheleine}, once more
  \item \emph{Say it:} a sentence of your own with \emph{die Wäscheleine}
  \item \emph{Recall:} say the German for cloth, then the German for tumble dryer, then say \emph{die Wäscheleine} again
\end{itemize}

\begin{tcolorbox}[breakable,colback=teal!4,colframe=teal!35!black,title={Before you move on}]
What does die Wäscheleine mean? (\textquotedblleft{}Clothes line\textquotedblright{}.) Say it once more.
\end{tcolorbox}
\section[die Wäscheklammer]{die Wäscheklammer — clothes peg}

\label{lesson:GE-C527-waescheklammer}



\begin{quote}
Before the new one: say the German for tumble dryer, then the German for clothes line.
\end{quote}

\subsection*{You'll want to know: die Wäscheklammer}
\textbf{die Wäscheklammer} — \textquotedblleft{}clothes peg\textquotedblright{}.

\begin{quote}
Für das große Laken braucht man viele Wäscheklammern.
\end{quote}

\begin{quote}
\emph{For the big sheet you need lots of pegs.}
\end{quote}

Plural: \emph{die Wäscheklammern}. \emph{Die Klammer} — clip; bracket.

\begin{cousinweb}
\emph{Wäscheklammer} is built on \emph{Wäsche} \textquotedblleft{}washing\textquotedblright{} and \emph{Klammer} \textquotedblleft{}clip\textquotedblright{}, inherited from Middle High German \emph{klammer} (EtymDB).
\end{cousinweb}

\begin{grammarlens}[title={What you've built}]
One more word for telling, explaining and describing at B1.
\end{grammarlens}

\subsection*{Your turn}
\begin{itemize}
\raggedright
  \item \emph{Say it:} \emph{die Wäscheklammer}
  \item \emph{Say it:} \emph{die Wäscheklammer}, once more
  \item \emph{Say it:} a sentence of your own with \emph{die Wäscheklammer}
  \item \emph{Recall:} say the German for tumble dryer, then the German for clothes line, then say \emph{die Wäscheklammer} again
\end{itemize}

\begin{tcolorbox}[breakable,colback=teal!4,colframe=teal!35!black,title={Before you move on}]
What does die Wäscheklammer mean? (\textquotedblleft{}Clothes peg\textquotedblright{}.) Say it once more.
\end{tcolorbox}
\section[bügeln]{bügeln — to iron}

\label{lesson:GE-C527-buegeln}



\begin{quote}
Before the new one: say the German for clothes line, then the German for clothes peg.
\end{quote}

\subsection*{You'll want to know: bügeln}
\textbf{bügeln} — \textquotedblleft{}to iron\textquotedblright{}.

\begin{quote}
Am Sonntagabend bügelt sie ihre Blusen für die Woche.
\end{quote}

\begin{quote}
\emph{On Sunday evening she irons her blouses for the week.}
\end{quote}

Regular: \emph{hat gebügelt}. In Switzerland also \emph{glätten}.

\begin{cousinweb}
\emph{bügeln} is built on \emph{Bügel} \textquotedblleft{}a curved handle, a hanger\textquotedblright{} (EtymDB); the link is usually explained by the bent handle of early irons.
\end{cousinweb}

\begin{grammarlens}[title={What you've built}]
One more word for telling, explaining and describing at B1.
\end{grammarlens}

\subsection*{Your turn}
\begin{itemize}
\raggedright
  \item \emph{Say it:} \emph{bügeln}
  \item \emph{Say it:} \emph{bügeln}, once more
  \item \emph{Say it:} a sentence of your own with \emph{bügeln}
  \item \emph{Recall:} say the German for clothes line, then the German for clothes peg, then say \emph{bügeln} again
\end{itemize}

\begin{tcolorbox}[breakable,colback=teal!4,colframe=teal!35!black,title={Before you move on}]
What does bügeln mean? (\textquotedblleft{}To iron\textquotedblright{}.) Say it once more.
\end{tcolorbox}
\section[falten]{falten — to fold}

\label{lesson:GE-C527-falten}



\begin{quote}
Before the new one: say the German for clothes peg, then the German for to iron.
\end{quote}

\subsection*{You'll want to know: falten}
\textbf{falten} — \textquotedblleft{}to fold\textquotedblright{}.

\begin{quote}
Nach dem Bügeln faltet er die Hemden.
\end{quote}

\begin{quote}
\emph{After ironing he folds the shirts.}
\end{quote}

Regular: \emph{hat gefaltet}. \emph{Die Hände falten} — to fold your hands.

\begin{cousinweb}
\emph{falten} is inherited from Old High German \emph{faltan}, from Proto-Germanic \emph{falþaną} \textquotedblleft{}to fold\textquotedblright{} (EtymDB); English \emph{fold} is a cousin.
\end{cousinweb}

\begin{grammarlens}[title={What you've built}]
One more word for telling, explaining and describing at B1.
\end{grammarlens}

\subsection*{Your turn}
\begin{itemize}
\raggedright
  \item \emph{Say it:} \emph{falten}
  \item \emph{Say it:} \emph{falten}, once more
  \item \emph{Say it:} a sentence of your own with \emph{falten}
  \item \emph{Recall:} say the German for clothes peg, then the German for to iron, then say \emph{falten} again
\end{itemize}

\begin{tcolorbox}[breakable,colback=teal!4,colframe=teal!35!black,title={Before you move on}]
What does falten mean? (\textquotedblleft{}To fold\textquotedblright{}.) Say it once more.
\end{tcolorbox}
